%==============================================================================
% JUAN DAVID UZHCA CAJAMARCA | SOFTWARE ENGINEER
% Master Professional CV | 100% ATS-Compliant (1 Page Strict)
% Compilador: pdfLaTeX / XeLaTeX
%==============================================================================

\documentclass[9.5pt, letterpaper]{article}

% --- Paquetes Fundamentales ---
\usepackage[utf8]{inputenc}
\usepackage[T1]{fontenc}
\usepackage{lmodern}
\usepackage[spanish, es-tabla]{babel}
\usepackage{geometry}
\usepackage{xcolor}
\usepackage{hyperref}
\usepackage{enumitem}
\usepackage{titlesec}
\usepackage{tabularx}
\usepackage{microtype}

% --- Geometría y Márgenes Calibrados para 1 Página Completa ---
\geometry{
    top=0.45in,
    bottom=0.45in,
    left=0.55in,
    right=0.55in
}

\setlength{\parindent}{0pt}
\setlength{\parskip}{0pt}
\pagestyle{empty}
\renewcommand{\familydefault}{\sfdefault}

% --- Paleta de Colores de Alta Fidelidad ---
\definecolor{navy}{HTML}{0F2A4A}       % Azul institucional oscuro
\definecolor{blueaccent}{HTML}{1D4ED8} % Azul técnico para enlaces y reglas
\definecolor{darktext}{HTML}{1E293B}   % Gris grafito oscuro (legibilidad óptima)
\definecolor{muted}{HTML}{475569}      % Gris medio para metadatos/subtítulos

\color{darktext}

% --- Configuración de Hipervínculos ---
\hypersetup{
    colorlinks=true,
    urlcolor=blueaccent,
    linkcolor=navy,
    pdftitle={CV - Juan David Uzhca - Software Engineer},
    pdfauthor={Juan David Uzhca Cajamarca}
}

% --- Títulos de Sección (Sin cajas, compatible con OCR/ATS) ---
\titleformat{\section}
    {\fontsize{10.5}{13}\selectfont\bfseries\color{navy}\uppercase}
    {}
    {0em}
    {}
    [\vspace{-2pt}{\color{blueaccent}\titlerule[1pt]}]

\titlespacing*{\section}{0pt}{4pt}{2.5pt}

% --- Listas Compactas sin sangría excesiva ---
\setlist[itemize]{
    noitemsep,
    topsep=1pt,
    partopsep=0pt,
    parsep=0.5pt,
    leftmargin=10pt,
    label={\color{blueaccent}\small$\bullet$}
}

\begin{document}

% --- ENCABEZADO MINIMALISTA Y ELEGANTE (Apto para escaneo ATS e impresión B/N) ---
\begin{center}
    {\fontsize{19}{22}\selectfont \textbf{\color{navy} JUAN DAVID UZHCA CAJAMARCA}}\\[2pt]
    {\fontsize{10.5}{12}\selectfont \textbf{Ingeniero en Ciencias de la Computación \textbar\ Software Engineer}}\\[2.5pt]
    \footnotesize
    Cuenca, Ecuador \textbar\ Open to Remote \textbar\ 
    \textbf{(+593) 99 851 8609} \textbar\ 
    \href{mailto:juandaviduzhca@hotmail.com}{juandaviduzhca@hotmail.com} \textbar\ 
    \href{https://linkedin.com/in/david-uzhca-5b5439442}{linkedin.com/in/david-uzhca-5b5439442} \textbar\ 
    \href{https://github.com/David5Uzhca}{github.com/David5Uzhca}
\end{center}

\vspace{-4pt}

% --- PERFIL PROFESIONAL ---
\section{Perfil Profesional}
\footnotesize
Ingeniero en Ciencias de la Computación enfocado en desarrollo \textbf{Full-Stack}, visión por computadora de alto rendimiento y arquitecturas concurrentes. Experiencia demostrada en el diseño e implementación de sistemas empresariales escalables (NestJS, Ionic/Angular) integrados a infraestructuras legacy (\textbf{IBM DB2 / AS400}), pipelines de visión acelerados por hardware (\textbf{C++17, OpenCV, CUDA 3.18x}) y agentes autónomos asistidos por IA (Google Gemini, FastAPI). Orientado a código limpio, modularidad y métricas medibles de rendimiento.

% --- HABILIDADES TÉCNICAS ---
\section{Habilidades Técnicas}
\footnotesize
\textbf{Lenguajes Core:} TypeScript, Python, C++17, Java (Spring Boot), SQL, C\#, Dart, JavaScript, HTML5/SCSS.\\
\textbf{Backend \& Arquitectura:} NestJS, FastAPI, Spring Boot, RESTful APIs, GraphQL (Strawberry), Microservicios, Node.js.\\
\textbf{Frontend \& Móvil:} Angular, Ionic Framework, React 18, Vite, Flutter, TailwindCSS, SCSS, RxJS.\\
\textbf{Visión por Computador \& IA:} OpenCV 4.x, CUDA / PyCUDA, PyTorch (Keypoint R-CNN), HOG + SVM, Google Gemini (Function Calling).\\
\textbf{Sistemas \& Concurrencia:} MPI (mpi4py), SPMD, Multithreading, Asynchronous I/O, IPC, Linux (Ubuntu/Kubuntu).\\
\textbf{Bases de Datos \& DevOps:} PostgreSQL, IBM DB2, AS400, Supabase, Git, Docker, Docker Compose, CMake, Postman.

% --- EXPERIENCIA Y PROYECTOS DESTACADOS ---
\section{Experiencia y Proyectos Destacados}

% PROYECTO 1: TESIS CENTROSUR
\noindent
\textbf{Sistema Empresarial de Reclamos Médicos} \textbar\ \textit{Lead Architect \& Full-Stack Developer} \hfill \textbf{2025 -- 2026}\\
\textit{Empresa Eléctrica Regional Centro Sur C.A. / Universidad Politécnica Salesiana} \hfill \textit{Cuenca, Ecuador}\\
\footnotesize\textbf{Stack:} NestJS, TypeScript, Ionic, Angular, IBM DB2, AS400, PostgreSQL, JWT, Chatbot IA Pepeluz
\begin{itemize}
    \item Diseñó e implementó una solución empresarial integral para la gestión y auditoría de reclamos médicos, reduciendo los tiempos de consulta administrativa mediante flujos automatizados de extremo a extremo.
    \item Desarrolló un backend modular desacoplado en \textbf{NestJS}, implementando autenticación JWT, RBAC y conectores optimizados hacia bases de datos transaccionales de nómina \textbf{IBM DB2 / AS400} y PostgreSQL.
    \item Construyó aplicaciones web y móviles multiplataforma en \textbf{Ionic / Angular} e integró el chatbot inteligente \textbf{"Pepeluz"} asistido por IA para consultas dinámicas de pólizas y expedientes en tiempo real.
\end{itemize}

\vspace{2.5pt}

% PROYECTO 2: SECURITY MONITOR
\noindent
\textbf{Security Monitor: Plataforma Híbrida de Vigilancia} \textbar\ \textit{Systems \& Computer Vision Engineer} \hfill \href{https://github.com/David5Uzhca/Security_Monitor_Based_on_Computer_Vision}{\textbf{[GitHub]}}\\
\textit{C++17, OpenCV 4.x, Hard Negative Mining, PyTorch (Keypoint R-CNN), CMake, Flask, Telegram Bot API} \hfill \textbf{2025 -- 2026}
\begin{itemize}
    \item Diseñó una arquitectura híbrida desacoplando la detección en tiempo real a alta velocidad en \textbf{C++17} del análisis semántico en \textbf{Python}, optimizando el descriptor \textbf{HOG + Linear SVM} con \textit{Hard Negative Mining} para mitigar falsos positivos.
    \item Minimizó la huella de memoria residente (RSS) mediante monitoreo nativo en el kernel Linux e implementó un buffer circular con comunicación IPC asíncrona hacia \textbf{PyTorch} para detección de caídas y despacho de alertas vía Telegram Bot.
\end{itemize}

\vspace{2.5pt}

% PROYECTO 3: UPS-GLAM
\noindent
\textbf{UPS-Glam: Plataforma Heterogénea con Aceleración GPU} \textbar\ \textit{Software Architect \& GPU Developer} \hfill \href{https://github.com/David5Uzhca/UPS-GLAM}{\textbf{[GitHub]}}\\
\textit{CUDA, PyCUDA, Python (FastAPI), Java 17 (Spring Boot), Flutter, Docker Compose, Supabase} \hfill \textbf{2025 -- 2026}
\begin{itemize}
    \item Implementó pipelines de convolución 2D de alta resolución acelerados mediante \textbf{kernels CUDA personalizados} en \textbf{PyCUDA}, optimizando transferencias host-to-device y alcanzando un incremento de velocidad de procesamiento de hasta \textbf{3.18x a 12x}.
    \item Arquitectó un microservicio en \textbf{FastAPI} con contexto CUDA dedicado en thread pool uniproceso para evitar fallos de concurrencia, integrando un backend transaccional en \textbf{Spring Boot} y una app cliente en \textbf{Flutter}.
\end{itemize}

\vspace{2.5pt}

% PROYECTO 4: NOVASHOP
\noindent
\textbf{NovaShop: Asistente Conversacional Multimodal con IA} \textbar\ \textit{Full-Stack \& GenAI Engineer} \hfill \href{https://github.com/David5Uzhca/ai_shopping_assistant}{\textbf{[GitHub]}}\\
\textit{Google Gemini API (Function Calling), FastAPI, GraphQL (Strawberry), React 18, TypeScript, ElevenLabs} \hfill \textbf{2025 -- 2026}
\begin{itemize}
    \item Construyó un agente autónomo de comercio electrónico usando \textbf{Google Gemini} con \textit{Tool / Function Calling} estricto para búsqueda de inventario y gestión de carritos sin alucinaciones, integrando síntesis de voz fotorrealista con \textbf{ElevenLabs}.
    \item Desarrolló una API Gateway híbrida en \textbf{FastAPI} con esquemas tipados \textbf{GraphQL} y persistencia transaccional en PostgreSQL.
\end{itemize}

\vspace{2.5pt}

% PROYECTO 5: SISTEMA DISTRIBUIDO AMBIENTAL
\noindent
\textbf{Sistema Distribuido de Monitoreo Ambiental} \textbar\ \textit{Concurrent Computing Engineer} \hfill \href{https://github.com/David5Uzhca/Sistema-de-Monitoreo-Ambiental-Concurrente}{\textbf{[GitHub]}}\\
\textit{Computación Paralela (MPI / mpi4py), Modelo SPMD, Python, PyQt6, Benchmarking} \hfill \textbf{2025}
\begin{itemize}
    \item Implementó una simulación paralela en clúster bajo el modelo \textbf{SPMD} con \textbf{mpi4py}, empleando operaciones colectivas (\textit{Scatter, Broadcast, Reduce, Barrier}) y evaluando la aceleración (\textit{Speedup}) y eficiencia conforme a la Ley de Amdahl con telemetría en \textbf{PyQt6}.
\end{itemize}

% --- EDUCACIÓN Y CERTIFICACIONES ---
\section{Educación y Certificaciones}
\noindent
\textbf{Ingeniería en Ciencias de la Computación} \textbar\ Universidad Politécnica Salesiana \hfill \textit{Cuenca, Ecuador \textbar\ Egreso: 2026}\\
\vspace{-2pt}
\begin{itemize}
    \item \textbf{Build, Train, and Deploy ML Models with Keras on Google Cloud} \textbar\ \textit{Google Cloud / Google Skills Boost} \hfill \textbf{Dic. 2025}\\
    Credencial oficial en arquitectura, entrenamiento y despliegue de modelos de Machine Learning en infraestructura cloud empresarial.
    \item \textbf{Certificate of Recognition: Graduation Project in English} \textbar\ \textit{Universidad Politécnica Salesiana} \hfill \textbf{Ago. 2026}\\
    Distinción académica otorgada por la defensa y sustentación del Proyecto de Titulación íntegramente en idioma inglés.
\end{itemize}

\end{document}
